% Methods: spatial blocking, accuracy assessment (H2) and conformal intervals.
% Design: D-031, D-069; PREREG sections 2 and 5-8, Amendments 2, 5 and 11.
% Every number comes from \metric{...}.

\subsection*{Spatial blocking}
The domain was tiled with \metric{params.block_size_km}~km squares aligned to the equal-area grid; each block takes the sector that contains its centre, which leaves \metric{blocks.n_total} blocks. They were assigned at random to training, calibration and test in proportions \metric{params.split_train}, \metric{params.split_calib} and \metric{params.split_test}, stratified by sector (seed \metric{params.seed}): \metric{blocks.n.train}, \metric{blocks.n.calib} and \metric{blocks.n.test} blocks (Supplementary Fig.~\ref{fig:studyarea}). The assignment was fixed before any evaluation, and a point or cell takes the split of its block. Every fitted quantity comes from training blocks only: the setup model and its threshold, the static-contour rule, the lidar constant, the grid offsets, the conformal scale model and the shape offsets of the hindcast audit. Conformal quantiles come from calibration blocks, and every reported accuracy and coverage from test blocks.

\emph{Block size.} Before any test-block value was read, an exponential model was fitted to the semivariogram of the topography-minus-lidar residuals in training blocks (\metric{params.vario_bin_km}~km lag bins up to \metric{params.vario_max_lag_km}~km, from at most \metric{params.vario_max_pairs} random point pairs). On synthetic white noise such a fit can return a long range, so its range is used only if spatial structure is detected (pre-registration Amendment~11). The statistic $T$ is the mean semivariance of pairs closer than \metric{params.vario_short_lag_km}~km over that of all pairs up to \metric{params.vario_max_lag_km}~km, with $p = (1 + \#\{T_{\mathrm{perm}} \le T\})/(N + 1)$ from $N = \metric{params.structure_perm_n}$ permutations of the residuals over the locations; structure is detected at $p \le \metric{params.alpha}$. If it is, and the practical range exceeds \metric{params.block_size_km}~km, the blocks grow to the next multiple of \metric{params.block_size_km}~km at or above the range and are reassigned with the same seed; with two lidar products the larger size applies. Blocks of \metric{params.block_size_sens_lo_km} and \metric{params.block_size_sens_hi_km}~km are a robustness check.

\subsection*{Accuracy of the topography}
\emph{Reference heights.} The primary reference is ICESat-2 ATL06 and ATL08 (Release~007)\todo{cite: ICESat-2 ATL06 and ATL08 Release 007 data sets (NSIDC DOIs) and Markus et al. 2017, Remote Sensing of Environment}. Segments were kept if their height and geoid height were finite, they came from a strong beam and their height lay within \metric{params.gross_error_m}~m of the product's own reference elevation; ATL06 segments also needed a zero quality summary, and ATL08 segments no multiple scattering and no cloud. The ATL08 water mask was not used: it is static, and would remove the very seabed under test. A point was used only if the seabed was dry when measured, that is, more than \metric{params.emergence_lag_days}~days after the emergence date of every topography cell within \metric{params.dry_radius_px_atl06} cell (ATL06) or \metric{params.dry_radius_px_atl08} cells (ATL08) of it, and the test is repeated with \metric{params.emergence_lag_sens_days}~days as a robustness check. ATL06 and ATL08 largely share photons, so they are evaluated separately and never pooled. GEDI L2A footprints\todo{cite: GEDI L2A Version 2 (Dubayah et al. 2021, doi:10.5067/GEDI/GEDI02\_A.002)} are an additional reference, and the substitute for ICESat-2, under the rules of pre-registration Amendment~2; GLO-30 heights over cells dry throughout its acquisition period are a further check.

\emph{Metrics and decision rule.} Errors are model minus reference. For the topography and each grid we report the bias, the root-mean-square error (RMSE), the normalised median absolute deviation (NMAD; \metric{params.nmad_scale} times the median absolute deviation from the median) and the \metric{params.err_pct_lo}th and \metric{params.err_pct_hi}th percentiles, pooled over test-block points where the reference and every model have a value, overall and per sector; the median of per-block values is a sensitivity. Each grid enters with one constant offset estimated on training blocks, as for the lidar constant, so that no grid is penalised for an undocumented datum; the grids as distributed are reported beside them. H2 is supported if the topography's test-block NMAD and RMSE are both lower than every grid's, with \metric{params.ci_pct}\% block-bootstrap intervals of the differences (\metric{params.bootstrap_n} resamples of test blocks, points pooled) that exclude zero, for both ATL06 and ATL08; it is partially supported if this holds for one. Per-sector intervals in sectors with few test blocks are descriptive.

\subsection*{Conformal intervals}
Height intervals come from split conformal prediction with locally normalised scores\citep{lei2018}. The score of a point is $|y - \hat{y}|/s(x)$, with $y$ the reference height on the DAHITI reference, $\hat{y}$ the topography and
\begin{equation}
  s(x) = \exp\!\left(\beta_0 + \beta_1 \log(1 + d) + \beta_2\,\kappa + \beta_3\,g\right),
\end{equation}
where $d$ is the interpolation distance, $\kappa$ the density of contour points (per \kmsq{} within \metric{params.contour_density_radius_m}~m) and $g$ the slope by central differences on the \metric{params.dem_res_m}~m grid. $\beta$ was fitted by least squares of $\log|y - \hat{y}|$ (residuals floored at \metric{params.abs_resid_floor_mm}~mm) on training blocks. On calibration blocks, $\hat{q}$ is the $\lceil (n+1)(1-\alpha) \rceil$-th smallest of the $n$ scores for nominal coverage $1-\alpha = \metric{params.conformal_level}$, and the interval is $\hat{y} \pm \hat{q}\,s(x)$. Coverage holds if calibration and test points are exchangeable; we assume exchangeability between blocks, which the disjoint calibration blocks are meant to make plausible. To check conditional coverage, $\hat{q}$ was computed per sector (Mondrian calibration)\todo{cite: Mondrian conformal prediction (Vovk et al., Algorithmic Learning in a Random World)} where a sector has at least \metric{params.mondrian_min_n} usable calibration points; other sectors use the pooled $\hat{q}$, and the value, count and fallback are reported for every sector. A point is usable if its reference, prediction and scale are all present; points outside the triangulation are excluded from calibration and coverage alike, and counted. Coverage on test blocks is reported overall and per sector with its Clopper--Pearson \metric{params.ci_pct}\% interval\todo{cite: Clopper \& Pearson 1934, Biometrika}, and counts as consistent with the nominal level if \metric{params.conformal_level} lies inside it. Points within a block are correlated, which makes that interval optimistic, so a block-bootstrap interval is also reported.
